Return-conditioned sequence policies turn desired return into an inference-time control input. This paper audits a failure mode in which larger candidate pools make a Decision-Transformer-style policy look more obedient to an ambitious return prompt while the selected trajectory becomes worse under the real objective. The mechanism is support mismatch: when the requested return lies beyond the offline behavior distribution and the proxy return underprices risky actions, top-score selection moves probability mass toward an extrapolative tail. We introduce a Return-Support Audit with three components: a tie-aware finite-pool accounting identity, a Return Fantasy Microscope for prompt, support, and real-utility gaps, and a Tail Phase Diagram that marks when candidate-count pressure helps, saturates, or hurts. In the out-of-support condition, naive top-score selection raises selected proxy return from $9.299$ to $9.468$ while real utility drops from $7.825$ to $7.381$ as candidate count grows from $1$ to $64$. The accounting identity matches Monte Carlo with maximum absolute error $0.00535$. In-support and anti-aligned controls behave as predicted, and likelihood or pilot-label repairs reduce the failure in this controlled environment. The result is a scoped audit protocol for deciding when high return prompts are evidence of improved utility and when they are extrapolation warnings.
